\appendix
% \chapter{Appendices} \label{appendix}

% In this appendix, I discuss two issues: The nature of so-called \is{bound marker}\textit{bound markers}, suffixes which might appear with relational modifiers in Japanese and have influence on their syntactic behavior, and the word class status of \textit{no}-modifiers.

\chapter{Bound markers and expanded modifiers}
% \addcontentsline{toc}{chapter}{Appendix A}
% \markboth{APPENDIX A}{Bound markers and expanded modifiers}

This section gives an overview of \textit{bound markers} -- also called \textit{relational nouns} -- those suffixes that form so-called \textit{expanded modifiers} out of Japanese relational modifiers discussed in Section \ref{RAJapanese} (cf. \citealt{nagano:2016:are, shimada:2018:relational}). First of all, bound markers specify the meaning of a \is{relational modifier}relational modifier. In the DP below, the bound marker -\textit{iri} `put in(-side)’ is added to the noun/relational modifier \textit{chīzu} `cheese’ and specifies the relationship between \textit{cheese} and the modified noun \textit{pan} `bread’. The expanded modifier formed out of this combination denotes precisely \textit{bread with cheese put inside}, whereas the DP \textit{chīzu-no pan} can mean anything from `bread made of cheese', `bread containing cheese' to `bread made for cheese’. Like relational modifiers, expanded modifiers appear exclusively with \textsc{no} in attributive position.

\setcounter{ExNo}{0}

\exg. chīzu-iri-no pan\\
cheese-put.in-\textsc{mod} bread\\
`bread with cheese put inside' \label{chizuiri}

The array of bound markers roughly corresponds to the list of categories for relational modifiers presented in Section \ref{RAJapanese}, but includes some more categories such as \textsc{taste}, \textsc{belonging}, \textsc{state}, \textsc{purpose}, \textsc{level} and \textsc{color}. The following is a (possible non-exhaustive) list, also taken from \citet[56–57]{nagano:2016:are}.\footnote{That \textit{no}-modifiers are often formed, or often contain, elements for word formation, preferably suffixes, is also noted in \citet[195–205]{muraki:2012:nihongo} who gives, besides his list of non-predicative modifiers, also a list of possible prefixes or suffixes often co-occurring with \textit{no}-modifiers.}

\ex. \emph{Nagano's list of bound markers}\\
\textsc{material}: \textit{sei} ‘made by’, \textit{iri} ‘added’\\
\textsc{origin}: \textit{sei} ‘made in’, \textit{kei} ‘descended from’, \textit{shushin} ‘coming from’, \textit{umare} ‘born in’\\
\textsc{shape}/\textsc{size}: \textit{kei} ‘shape’, \textit{kata/gata} ‘shape, size’\\
\textsc{taste}: \textit{aji/mi} ‘taste’, \textit{fūmi} ‘taste’\\
\textsc{type}: \textit{sei} ‘type, nature’, \textit{gata} ‘type’, \textit{kei} ‘type’\\
\textsc{state}: \textit{jō} ‘state’, \textit{jōtai} ‘state’, \textit{sugata} ‘wearing’\\
\textsc{belonging}: \textit{kumi} ‘group’, \textit{ha} ‘group, school’, \textit{shugi} ‘-ism’, \textit{shozoku} ‘belonging to’\\
\textsc{similarity}: \textit{fū} ‘like’, \textit{ryū} ‘like, in the style of’\\
\textsc{possession}/\textsc{ingredient}: \textit{tsuki} ‘with’, \textit{mochi} ‘with’, \textit{iri} ‘added’, \textit{ari} ‘with’, \textit{nashi} ‘-less, without’, \textit{darake} ‘-ful, -ridden’\\
\textsc{purpose}/\textsc{target}: \textit{yō} ‘for’, \textit{muke} ‘meant for’, \textit{senyō} ‘exclusively for’\\
\textsc{location}: \textit{mae}/\textit{zen} ‘front’, \textit{shita}/\textit{ka} ‘under’, \textit{naka} ‘in, inside’, \textit{ue}/\textit{jō} ‘on, above’, \textit{chū} ‘inside’, \textit{kan} ‘between’, \textit{iki} ‘bound for’, \textit{hatsu} ‘departing from’, \textit{muki} ‘toward, faced to’, \textit{kake} ‘hanged on’\\
\textsc{time}: \textit{mae} ‘before’, \textit{go} ‘after’, \textit{chū} ‘during’\\
\textsc{status}/\textsc{possession}: \textit{jin} ‘nationality’, \textit{shi} ‘specialist’, \textit{fu} ‘female’, \textit{kan} ‘official’, \textit{ko} ‘worker’, \textit{jo} ‘female’, \textit{toshite} ‘as’\\
\textsc{level}: \textit{kyū} ‘level’, \textit{reberu} ‘level’, \textit{do} ‘degree’, \textit{i} ‘level’\\
\textsc{topic}: \textit{jō} ‘on, about’, \textit{garami} ‘about, related to’, \textit{kanren} ‘related to’\\
\textsc{color}: \textit{shoku} ‘color’, \textit{iro} ‘color’

Some of these \is{bound marker}markers are derived from nouns which can be used on their own, such as \textit{jōtai} `state’, while other bound markers cannot be used independently, for example -\textit{darake} `-ful, -ridden’.

These markers come with interesting syntactic peculiarities. When they are added to a, formerly non-predicative, \is{relational modifier}relational modifier, the resulting expanded modifier becomes possible in predicative position as shown below \citep{nagano:2016:are, shimada:2018:relational}.\footnote{The following data points were also part of the interview conducted for relational modifiers and reported on in Chapter \ref{direct japanese}.}

\exg. Kono pan-wa chīzu-iri-d-a.\\
this bread-\textsc{top} cheese-put.in-\textsc{pred-cop.prs}\\
`This bread is with cheese inside.’

\exg. Kono kuruma-wa Nihon-sei-d-a.\\
this car-\textsc{top} Japan-produced-\textsc{pred-cop.prs}\\
`This car is Japanese. (= was produced in Japan)'

\largerpage
Following the analysis of \textsc{de} used for coordination as the predicative copula (cf. Section \ref{copulaelements}), we expect \is{bound marker}expanded modifiers to be also able in coordinative structures. This is indeed the case. \is{coordination} \is{coordinative copula}

\ex. \ag. chīzu-iri-de herushi-na pan\\
cheese-put.inside-\textsc{pred} healthy-\textsc{mod} bread\\
`bread that is cheesy (= cheese put inside) and healthy'
\bg. Nihon-sei-de jōbu-na kuruma\\
Japan-produced-\textsc{pred} robust-\textsc{mod} car\\
`car that is Japanese (= produced in Japan) and robust'

Finally, while paraphrasing \textsc{no} with the canonical copula \textsc{de-ar-u} was never possible with \is{relational modifier}relational modifiers, this operation is possible with expanded modifiers.

\exg. Nihon-sei-de-ar-u tokei\\
Japan-produced-\textsc{pred-cop-prs} watch\\
`a Japanese watch' \null \hfill \citep[75]{shimada:2018:relational}

On the other hand, in contrast to canonical relational modifiers, these extended modifiers do not exhibit nominal use.

\exg. *Chūgoku-sei-ga kawat-ta\\
China-produced-\textsc{nom} change-\textsc{pst}\\
(`(produced in) China changed’)

This suggests that they either do not have a nominal status anymore, or that they have become lexicalized. Adjectival status, at first glance, seems equally doubtful. They sill remain \is{gradability}non-gradable, in combination with the relative degree adverb \textit{totemo} `very’ and in comparative clauses \ref{expandedgrada}, and cannot be \is{nominalization}nominalized via the suffix -\textit{sa} `ness’ \ref{expandedsa}.

\ex. \label{expandedgrada} \ag. *totemo Nihon-san-no kuruma\\
very Japanese-production-\textsc{mod} car\\
(`very Japanese car’)
\bg. *Watashi-ga okusan-no kuruma-to kurabe-reba Nihon-san-no kuruma-wo kat-ta.\\
I-\textsc{nom} wife-\textsc{mod} car-\textsc{com} compare-if Japan-production-\textsc{mod} car-\textsc{acc} buy-\textsc{pst}\\
(`Compared to my wife, I bought a more Japanese car.’)

\ex. \label{expandedsa} \ag. *Chūgoku-sei-sa\\
China-produce-ness\\
\bg. *Nihon-san-sa\\
Japan-produce-ness\\

However, at least for some suffixes, especially those derived from lexical verbs, \is{gradability}gradability in comparative constructions with the particle \textit{yori} `than’ is possible. Example \ref{lawson}, in which the quantity of cheese is compared, was judged acceptable by three out of five consultants.

\exg. LAWSON-no pan-to kurabe-reba motto chīzu-iri-no pan-wo kat-ta.\\
LAWSON-\textsc{mod} bread-\textsc{com} compare-if more cheese-put.inside-\textsc{mod} bread-\textsc{acc} buy-\textsc{pst}\\
`I bought bread with more cheese inside compared to bread from LAWSON.’ \label{lawson}

That \textit{chīzu-iri} is inherently verbal is visible when paraphrasing this modifier with a verbal relative clause, either intransitive or transitive. In these constructions, comparison is equally possible.

\ex. \ag. LAWSON-no pan-yori chīzu-ga hait-te-i-ru pan\\
LAWSON-\textsc{mod} bread-than cheese-\textsc{nom} be.inside-\textsc{ger}-be-\textsc{prs} bread\\
`bread with more cheese inside than bread from LAWSON’
\bg. LAWSON-no pan-yori chīzu-wo ire-ta pan\\
LAWSON-\textsc{mod} bread-than cheese-\textsc{acc} put.inside-\textsc{pst} bread\\
`bread in which X has put more cheese than in bread from LAWSON’

In order to give a syntactic explanation for predicative use, \citet{shimada:2018:relational} adhere to a \is{bound marker}KIND-head following \citet{adger:2013:syntax}. This KIND-head, which is more of a variable, is present in predicative position and is what is really modified by the \is{relational modifier}relational modifier, meaning that essentially we are dealing with a modification structure (see especially \citealt[§ 4]{adger:2013:syntax} for the underlying analysis of relational nouns). This proposal is in spirit very similar to the semantic \textit{kind-reading} in \citet{mcnally:2004:relational}. Recall the following example from Catalan.  \is{Catalan}

\exg. La tuberculosi pot ser pulmonar. (Catalan)\\
the tuberculosis can be pulmonary\\
`The tuberculosis could be (of the) pulmonary (kind/type).’ \\ \null \hfill \citep[189]{mcnally:2004:relational}

Also note the following example in \is{English} English from \citet{bauer:2013:oxford} containing the \is{relational modifier}relational adjective \textit{nuclear} in predicative position. Predicative use is further ameliorated by a measure phrase containing an endpoint modifier (\citealt{fabregas:2014:adjectival}, see \citealt[80]{shimada:2018:relational}).

\ex. 75\% of French electricity is nuclear. \hfill \citep[318]{bauer:2013:oxford}

In this example, a noun such as \textit{kind} or \textit{type} seems to be understood. Therefore, the sentence can be paraphrased accordingly either with the overt noun repeated or a generic \textit{kind}-noun. \is{kind-reading}

\ex. \a. 75\% of French energy is nuclear energy.
\b. 75\% of French energy is of the nuclear kind/type.

Japanese seems to overtly express such a KIND-head in the form of \is{bound marker}bound markers. This essentially suggests the structure displayed in \figref{structure relational modifiers and bound markers}, based on \ref{chizuiri}. To be concrete, what really is selected by ModP is a DP whose lexical head is the bound marker and in whose extended projection the modifier \textit{chīzu} resides.

\begin{figure}
\centering
 \begin{tikzpicture}[baseline=0pt, remember picture]
\tikzset{level distance=25pt, sibling distance=5pt}
\Tree [.DP\textsubscript{internal} {} [.D′\textsubscript{internal} [.FP [.ModP {} [.Mod′ [.DP\textsubscript{internal} {} [.D′\textsubscript{internal} [.FP [.XP \edge[roof]; chīzu ] [.F′ [.NP \edge[roof]; iri ] F ] ] D ] ] [.Mod no ] ] ] [.F′ [.NP \edge[roof]; pan ] F ] ] D ] ]
\end{tikzpicture}
\caption{Structure of relational modifiers with bound markers (= \ref{chizuiri})}
\label{structure relational modifiers and bound markers}
\end{figure}

The only question is, whether \textit{chīzu} is not selected by \is{MOD@\textsc{mod}}\textsc{mod} or whether this head is covert again. Presumably, we see some kind of lexicalization here and therefore this syntactic head might not be necessary. However, a \is{relational modifier}direct counterargument against a status of \textit{chīzu-iri} as lexical compound comes from structures as in \ref{oishichizu}. \is{bound marker}

\exg. [[oishi-i chīzu-iri-no] pan]\\
tasty-\textsc{mod} cheese-put.in-\textsc{mod} bread\\
`bread with tasty cheese put inside' \label{oishichizu}

While the \textit{i}-adjective \textit{oishi-i} `tasty’ can modify \textit{pan} `bread’, it can also modify \textit{cheese} as indicated by the \is{bracketing ambiguity}brackets. If we were dealing with a lexical compound the first part should not be able to be modified.

There are thus many issues left open for the research on bound markers and expanded modifiers to be addressed in the future.

\chapter{A taxonomy of \textit{no}-modifiers}
% \addcontentsline{toc}{chapter}{Appendix B}
% \markboth{APPENDIX B}{A taxonomy of no-modifiers}

In this section, I present an overview over the taxonomy of \textit{no}-modifiers I established. More details can be found in the data paper contained in the repository under \url{https://osf.io/u9txp/}.

I have tested every lexeme in the data set for the following five variables. \is{word classes}

\ex. \a. occurrences in attributive use
\b. element in attributive position: \textsc{na} or \textsc{no}
\c. predicative use with copula
\d. referential use with case particles for subject and object
\e. nominalization

Based on these variables I have established the following taxonomy of \textit{no}-modifiers.

\ex. Taxonomy of \textit{no}-modifiers
\a. \textit{no}-adjectives (Group 1)
\b. Dual membership: Adjectival use with \textsc{no} and \textsc{na} but not as noun (Group 2)
\c. Dual membership: Use as noun and adjectival use as \textit{na}-adjective, in some instances as \textit{no}-modifiers (Group 3)
\d. Pure nouns (Group 4)
\e. Direct modifiers (Group 5)

There are several important things to note: Group 3 can be further divided into lexemes where attributive use with adjectival (qualitative) semantics is only possible with \textsc{na}, and those where it is also possible with \textsc{no}. An instance of the former would be \textit{kenkō} `healthy, health’ an instance of the latter \textit{byōdō} `equal, equality’. Since this is a semantic factor, I have not established a different group, but sorted these modifiers into two subgroups. Second, since modification status and word class are treated as orthogonal issues, for any given lexeme nominal or adjectival status is independent of its role as direct modifier. For some lexemes, one could argue that they should rather be put in another group than group 5. However, since direct modification is such an important part of this study, I have dedicated a whole group to direct modifiers. Finally, note that group 5 dan be subdivided into those direct modifiers without a nominal status (group 5a) and those with a nominal status (group 5b). \is{word classes}

Taking all these factors into account, an overview of the groups and relevant characteristics looks as \is{nominalization}displayed \is{MOD@\textsc{mod}}in \ref{taxonomy list}.

\ex. \label{taxonomy list} Taxonomy of \textit{no}-modifiers and their characteristics
\a. Group 1 (\textit{no}-adjectives): \textsc{mod} is realized as /no/, predicative use possible, referential use impossible, nominalization possible (\textit{mumei} `unknown')
\b. Group 2 (\textit{no}/\textit{na}-adjectives): \textsc{mod} is realized as /no/ or /na/, predicative use possible, referential use impossible, nominalization possible (\textit{tokubetsu} `special')
\c. Group 3a (\textit{na}-adjectives/nouns): \textsc{mod} is realized as /no/ or /na/, predicative use possible, referential use possible, nominalization marginally possible, qualitative modification only with /na/ (\textit{kenkō} `health(y)')
\d. Group 3b (\textit{no}/\textit{na}-adjectives/nouns): \textsc{mod} is realized as /no/ or /na/, predicative use possible, referential use possible, nominalization impossible, qualitative modification only with /na/ (\textit{byōdō} `equal(ity)')
\e. Group 4 (nouns): \textsc{mod} is realized as /no/ only, predicative use possible, referential use possible, nominalization impossible, qualitative modification marginally possible with /no/ (\textit{byōki} `sickness')
\e. Group 5a (direct \textit{no}-modifiers): \textsc{mod} is realized as /no/, predicative use impossible, referential use impossible, nominalization marginally possible (\textit{seirai} `innate')
\e. Group 5b (direct \textit{no}-modifiers with nominal status): \textsc{mod} is realized as /no/, predicative use impossible, referential use impossible, nominalization impossible (\textit{tagai} `each other')

%\begin{table}
%\begin{tabular}{llllllll}
%\lsptoprule
%\textbf{Group} & 1 & 2 & 3a & 3b & 4 & 5a & 5b\\
%\midrule
%Example & \shortstack{mumei\\`unknown’} 
 %       & \shortstack{tokubetsu\\`special’} 
  %      & \shortstack{kenkō\\`health(-y)’} 
   %     & \shortstack{byōdō\\`equal(-ity)’} 
    %    & \shortstack{byōki\\`sickness’} 
     %   & \shortstack{seirai\\`innate’} 
      %  & \shortstack{tagai\\`each other’} \\
%\textsc{mod} & no & no/na & no/na & no/na & no & no & no \\
%Predicative Use & \checkmark & \checkmark & \checkmark & \checkmark & \checkmark & $\times$ & %$\times$ \\
%Referential Use & $\times$ & $\times$ & \checkmark & \checkmark & \checkmark & $\times$ & %\checkmark \\
%Nominalization & \checkmark & \checkmark & marginal & $\times$ & $\times$ & marginal & $\times$ %\\
%\lspbottomrule
%\end{tabular}
%\caption{Taxonomy of \textit{no}-Modifiers}
%\label{tab:taxonomy-no-modifiers}
%\end{table}

The emerging picture shows a spectrum from more adjectival to nominal status. Group 1 fulfill the most adjectival criteria, while the other extreme, group 4, consists of true nouns that fulfill none of the criteria defined for adjectives.

\largerpage[2]
Below, tables with illustrative examples of each group alongside their occurrences in the relevant environments on the \is{BCCWJ}BCCWJ (analysis finalized in 08/2023; cf. \citealt{kohlich:2023:suppl} and Section \ref{corpusanalysis}) are given. Group 1 is displayed in \tabref{table group 1}, group 2 in \tabref{table group 2}, group 3 in \tabref{table group 3} and group 4 in \tabref{table group 4}. Examples of direct modifiers without nominal status are given in \is{word classes}\tabref{table group 5}.\clearpage

\begin{table}[p!]
\centering
\begin{tabularx}{\textwidth}{Xrrrrr}
\lsptoprule
\textbf{lexeme} & \textbf{-no} & \textbf{-na} & \textbf{copula} & \textbf{case} & \textbf{-sa} \\
\midrule
\textit{mumei} `unknown’    & 217 & 11 & 46  & 4  & 0 \\
\textit{hadashi} `barefoot’ & 101 & 0  & 58  & 8  & 0 \\
\textit{kogata} `small-size’& 522 & 8  & 125 & 10 & 0 \\
\lspbottomrule
\end{tabularx}
\caption{Group 1: \textit{no}-adjectives, 215 lexemes}
\label{table group 1}
\end{table}

\begin{table}[p!]
\centering
\begin{tabularx}{\textwidth}{Xrrrrr}
\lsptoprule
\textbf{lexeme} & \textbf{-no} & \textbf{-na} & \textbf{copula} & \textbf{case} & \textbf{-sa} \\
\midrule
\textit{pittari} `fitting perfectly’ & 301 & 69  & 338 & 0   & 0 \\
\textit{muyō} `unnecessary’          & 196 & 115 & 182 & 4   & 2 \\
\textit{kō’un} `good luck’           & 106 & 134 & 209 & 0   & 0 \\
\lspbottomrule
\end{tabularx}
\caption{Group 2, \textit{no/na}-modifiers, 80}
\label{table group 2}
\end{table}

\begin{table}[p!]
\centering
\begin{tabularx}{\textwidth}{Xrrrrr}
\lsptoprule
\textbf{lexeme} & \textbf{-no} & \textbf{-na} & \textbf{copula} & \textbf{case} & \textbf{-sa} \\
\midrule
\textit{kurashikku} `classic’        & 102 & 93  & 37  & 50  & 1 \\
\textit{kyū} `urgent, urgency’       & 21  & 605 & 138 & 118 & 1 \\
\textit{manzoku} `pleased’           & 218 & 138 & 345 & 248 & 0 \\
\lspbottomrule
\end{tabularx}
\caption{Group 3: adjectives/nouns, 105}
\label{table group 3}
\end{table}

\begin{table}[p!]
\centering
\begin{tabularx}{\textwidth}{Xrrrrr}
\lsptoprule
\textbf{lexeme} & \textbf{-no} & \textbf{-na} & \textbf{copula} & \textbf{case} & \textbf{-sa} \\
\midrule
\textit{hantai} `opposite, contrary’ & 1140 & 23 & 656 & 510 & 0 \\
\textit{gūzen} `accident(al)’       & 327  & 20 & 371 & 124 & 0 \\
\textit{erīto} `elite’               & 77   & 1  & 63  & 51  & 0 \\
\lspbottomrule
\end{tabularx}
\caption{Group 4, nouns: 147}
\label{table group 4}
\end{table}

\begin{table}[p!]
\centering
\begin{tabularx}{\textwidth}{Xrrrrr}
\lsptoprule
\textbf{lexeme} & \textbf{-no} & \textbf{-na} & \textbf{copula} & \textbf{case} & \textbf{-sa} \\
\midrule
\textit{kōzen} `publicly known’   & 40  & 0  & 0  & 0  & 0 \\
\textit{seirai} `innate’          & 92  & 0  & 0  & 5  & 0 \\
\textit{eien} `eternity, eternal’ & 793 & 12 & 38 & 31 & 0 \\
\lspbottomrule
\end{tabularx}
\caption{Group 5: direct modifiers, 511}
\label{table group 5}
\end{table}
\clearpage

Finally, subgroup 5b of \textit{direct nominals} is presented in \tabref{table group 5b}. A direct nominal is a direct modifier with at least 5 occurrences with case particles or a ratio higher than 5:1 for nominal use. This applies to a total of 90 lexemes. Good examples are most \is{relational modifier}relational modifiers, including modifiers of material such as \textit{tetsu} `iron’, but also \textit{ōkuchi} `large (amount)’ or \textit{shōsai} `detailed, details’. \is{word classes}


\begin{table}[t!]
\centering
\begin{tabularx}{\textwidth}{Xrrrrr}
\lsptoprule
\textbf{lexeme} & \textbf{-no} & \textbf{-na} & \textbf{copula} & \textbf{case} & \textbf{-sa} \\
\midrule
\textit{tetsu} `iron’               & 1117 & 2   & 39 & 348 & 0 \\
\textit{ōkuchi} `large (amount)’    & 64   & 1   & 2  & 47  & 0 \\
\textit{shōsai} `detailed, details’ & 32   & 787 & 46 & 472 & 2 \\
\lspbottomrule
\end{tabularx}
\caption{Group 5b: direct nominal modifiers, 90}
\label{table group 5b}
\end{table}

~




